\chapter{Integración}
\label{cap:integracion}

La integral de Riemann se construye mediante particiones y sumas superiores e inferiores. El capítulo reúne los criterios de integrabilidad y el teorema fundamental del cálculo \cite{apostol1974,spivak2008,rudin1976}.

\section{Particiones de un intervalo}
Particiones, refinamientos, malla y variación de una función en subintervalos.

\section{Sumas superiores e inferiores}
Sumas de Darboux y criterio de integrabilidad.

\section{Integral de Riemann}
Definición, existencia para funciones acotadas y ejemplos básicos.

\section{Criterios de integrabilidad}
Integrabilidad de funciones continuas y monótonas y criterio de Lebesgue para Riemann.

\section{Funciones de Riemann integrables}
Caracterizaciones y ejemplos de funciones acotadas que son integrables.

\section{Propiedades de la integral}
Linealidad, monotonía, aditividad por intervalos y estimaciones.

\section{Teorema fundamental del cálculo}
Relación entre derivación e integración y sus dos partes.

\section{Integración por sustitución}
Cambio de variable en integrales definidas e indefinidas.

\section{Integración por partes}
Fórmula de integración por partes y selección de factores adecuados.

\section{Integrales impropias}
Definición mediante límites y criterios básicos de convergencia.
